The displayed formulas are the chain rule on $\tau^e = \exp(e\log\tau)$. For the non-membership, fix $x = x_0$, write $e_0 := e(x_0)$, and suppose $\tau^{e_0}\log\tau = \sum_j g_j(\tau)\,\tau^{c_j}$ with $g_j$ smooth and $\{c_j\}$ an admissible exponent family evaluated at $x_0$: distinct reals, only finitely many below any level. Absorb the finitely many terms with $c_j = e_0$ into the left side and divide by $\tau^{e_0}$:
\[
\log\tau - g(\tau) \;=\; \sum_j g_j(\tau)\,\tau^{d_j}, \qquad d_j := c_j - e_0 \ \text{distinct and nonzero},
\]
with $g$ smooth. Split the right side into the positive part $P$ ($d_j > 0$: finitely many terms below level $1$ plus an $O(\tau)$ tail, so $P = o(1)$) and the negative part $R$ ($d_j < 0$: finitely many terms, by admissibility). Then $R = \log\tau + O(1)$, so in particular $R = o(\tau^{-\varepsilon})$ for every $\varepsilon > 0$. We raise the exponents of $R$ one at a time. Let $d_{\min} < 0$ be the least exponent present in $R$; every term of $\tau^{|d_{\min}|} R = \sum g_j \tau^{d_j + |d_{\min}|}$ has nonnegative exponent, so it extends continuously to $\tau = 0$ with limit $g_{j_0}(0)$, where $j_0$ is the (unique, exponents being distinct) index with $d_{j_0} = d_{\min}$. But $\tau^{|d_{\min}|} R = \tau^{|d_{\min}|}(\log\tau + O(1)) \to 0$; hence $g_{j_0}(0) = 0$, and $g_{j_0} \tau^{d_{\min}} = (g_{j_0}/\tau)\,\tau^{d_{\min}+1}$ with $g_{j_0}/\tau$ smooth: the term's exponent rises by one (merging with another term if the exponents now coincide). Each step raises the least exponent, and after finitely many steps every exponent in $R$ is nonnegative, making $R$ bounded near $0$. Then $\log\tau = g + P + R = O(1)$, contradicting $\log\tau \to -\infty$.
